\section{Objet Assembly}
\label{sec:objets_assembly}

L'\textbf{Assembly Object} (classe 0x04) est l'objet central de la messagerie implicite : il regroupe, dans un seul bloc de données contigu, plusieurs attributs qui seraient sinon dispersés dans plusieurs objets/instances. Cela évite de multiplier les connexions cycliques et permet un échange efficace en une seule trame.

\subsection{Les trois assemblies typiques}

\begin{center}
    \begin{tabular}{|p{2.8cm}|p{5cm}|p{6.5cm}|}
        \hline
        \textbf{Assembly} & \textbf{Sens du point de vue de l'Adapter} & \textbf{Contenu typique} \\
        \hline
        Input Assembly & Adapter $\rightarrow$ Scanner & États, mesures, retours de capteurs, codes de défaut. \\
        \hline
        Output Assembly & Scanner $\rightarrow$ Adapter & Commandes, consignes, ordres de marche. \\
        \hline
        Configuration Assembly & Scanner $\rightarrow$ Adapter (généralement à l'ouverture de connexion) & Paramètres de configuration transmis une fois à l'établissement de la connexion. \\
        \hline
    \end{tabular}
\end{center}

\begin{UPSTIinfor}{Point de vue Scanner vs point de vue Adapter}
    Attention au sens : ce que l'Adapter appelle son \emph{Input Assembly} correspond, côté Scanner (votre automate), à une donnée reçue en \textbf{entrée}. Ce que l'Adapter appelle son \emph{Output Assembly} correspond, côté Scanner, à une donnée envoyée en \textbf{sortie}. La documentation constructeur (fichier EDS) se place presque toujours du point de vue de l'Adapter -- gardez ce repère en tête au moment de câbler vos variables Codesys.
\end{UPSTIinfor}

\subsection{Dimensionnement}
Chaque Assembly a une \textbf{taille fixe}, exprimée en octets, définie soit par une instance standard proposée par l'équipement, soit par une configuration (\emph{assembly configurable}) lorsque l'équipement le permet. Ce dimensionnement doit être décidé \emph{avant} la mise en \oe uvre :
\begin{itemize}
    \item Lister précisément les données à échanger dans chaque sens.
    \item Choisir un format compact (bits groupés en octets/mots, éviter le gaspillage).
    \item Vérifier la correspondance entre la taille de l'Assembly déclarée côté Adapter et celle attendue côté configuration Scanner : une incohérence empêche l'établissement de la connexion.
\end{itemize}

\begin{UPSTIidee}{Lien avec votre code orienté objet}
    Une bonne pratique de conception consiste à faire correspondre un Assembly Object à une \textbf{structure de données} (ou à une classe simple) côté Codesys, représentant l'état échangé avec l'équipement distant. L'objet logiciel qui pilote l'équipement (ex. \texttt{CVariateur}, \texttt{CCameraVision}) encapsule alors la mise en forme/lecture de cette structure, et expose à son tour des propriétés et méthodes de haut niveau au reste de l'application -- la complexité du mapping CIP reste cachée à l'intérieur de l'objet.
\end{UPSTIidee}
